\metatitle{Real-rooted polynomials: open problems}

\metadescription{Open conjectures on real-rootedness and interlacing of
combinatorial polynomials, with definitions, formulas, first polynomials,
tested ranges and partial results, and a list of counterexamples.}


\section[realrootedOpenProblems]{Open problems on real-rootedness}

This page collects open conjectures about
\hyperref[realRootedPolynomial]{real-rootedness} and
\hyperref[interlacingPolynomials]{interlacing} of combinatorial polynomials.
Each entry gives the combinatorial meaning, a formula or recurrence when one
is known, the first few polynomials, the range in which the statement has
been tested, and what has been proved.  The interlacing convention
$f\interl g$ is the one on the
\hyperref[interlacesPolynomial]{interlacing page}.  A family
$(P_n)_n$ \emph{interlaces consecutively} if $P_{n}\interl P_{n+1}$ for
all $n$ in the stated range.

Unless a reference is given, the tested ranges come from our exact
computations: real-rootedness was decided with Sturm sequences on the
square-free parts, so repeated roots are allowed.  A finite check is
evidence, not a proof.  The last section lists
\hyperref[realRootedCounterexamples]{counterexamples}: families that look
real-rooted for small parameters but are not.

The general theory is on the pages on
\hyperref[polynomialRootIntro]{real-rooted polynomials} and
\hyperref[interlacingPolynomials]{interlacing polynomials}; proved examples
are on the pages on \hyperref[realRootedWords]{permutation statistics},
\hyperref[realRootedCatalan]{Catalan families},
\hyperref[sturmSetPartitions]{set partitions},
\hyperref[realRootedTableaux]{tableaux and polytopes} and
\hyperref[realRootedMatching]{graph polynomials}.


\section[openProblemsPermutations]{Permutation statistics}

\subsection[openProblemsCycleRestricted]{Permutations with long cycles}

For $k\geq 1$, let $\symS_n^{(k)}$ denote the permutations in $\symS_n$
all of whose cycles have length at least $k$.  Thus $\symS_n^{(1)}=\symS_n$
and $\symS_n^{(2)}$ is the set of derangements.  Define
\[
  E_{n,k}(t)\coloneqq\sum_{\pi\in\symS_n^{(k)}}t^{\exc(\pi)},
  \qquad
  D_{n,k}(t)\coloneqq\sum_{\pi\in\symS_n^{(k)}}t^{\des(\pi)}.
\]

\begin{conjecture}[Alexandersson, 2026]
\label{conjLongCycleExcedances}
For every $k\geq 3$, the polynomials $E_{n,k}(t)$, $n\geq k$, are
real-rooted and interlace consecutively.
\end{conjecture}

The excedances of the cycles of length $m$ are counted by $tA_{m-1}(t)$,
where $A_m(t)$ is the \hyperref[eulerianPolynomial]{Eulerian polynomial}
with $A_0=1$.  The exponential formula gives
\[
  \sum_{n\geq 0}E_{n,k}(t)\frac{z^n}{n!}
  =\exp\Bigl(\sum_{m\geq k}tA_{m-1}(t)\frac{z^m}{m!}\Bigr).
\]
Inserting $n$ into a cycle, or closing a new $k$-cycle through $n$, gives
\[
  E_{n,k}(t)=(n-1)t\,E_{n-1,k}(t)+t(1-t)E'_{n-1,k}(t)
  +\binom{n-1}{k-1}tA_{k-1}(t)\,E_{n-k,k}(t).
\]
Indeed, inserting $n$ after $i$ in its cycle keeps the excedance count if
$i$ was an excedance, and raises it by one otherwise.  Since every
$tA_{m-1}(t)$ is palindromic, so is every $E_{n,k}(t)$.  For $k=3$,
\[
\begin{array}{rl}
\toprule
n & E_{n,3}(t) \\
\midrule
3 & t+t^2 \\
4 & t+4t^2+t^3 \\
5 & t+11t^2+11t^3+t^4 \\
6 & t+36t^2+86t^3+36t^4+t^5 \\
7 & t+92t^2+477t^3+477t^4+92t^5+t^6 \\
8 & t+211t^2+2143t^3+4278t^4+2143t^5+211t^6+t^7 \\
\bottomrule
\end{array}
\]
Evidence: real-rootedness and consecutive interlacing hold for $k\leq 6$
and $n\leq 22$.  The case $k=2$ is the
\hyperref[derangementPolynomial]{derangement polynomial}, where the extra
term is $(n-1)tE_{n-2,2}$ and the recurrence gives an interlacing proof by
induction.  For $k\geq3$ the lagged term carries the extra factor
$A_{k-1}(t)$, and we know no proof.

\begin{conjecture}[Alexandersson, 2026]
\label{conjLongCycleDescents}
For every $k\geq 3$, the polynomials $D_{n,k}(t)$ are real-rooted.
\end{conjecture}

By the \name{Ira Gessel}--\name{Christophe Reutenauer} theorem
\cite{GesselReutenauer1993}, the quasisymmetric generating function for
descent sets over $\symS_n^{(k)}$ is the degree-$n$ part of
$\prod_{j\geq k}H[L_j]$, where $L_j$ is the Lie character and
$H=\sum_{m\geq0}\completeH_m$.  Counting necklaces gives
\[
  \frac{D_{n,k}(t)}{(1-t)^{n+1}}
  =\sum_{N\geq1}t^{N-1}\,[z^n]\prod_{j\geq k}(1-z^j)^{-M_j(N)},
  \qquad
  M_j(N)=\frac1j\sum_{d\mid j}\mu(d)N^{j/d},
\]
for $n\geq1$.  For $k=3$,
\[
\begin{array}{rl}
\toprule
n & D_{n,3}(t) \\
\midrule
3 & 2t \\
4 & 3t+3t^2 \\
5 & 6t+12t^2+6t^3 \\
6 & 12t+68t^2+68t^3+12t^4 \\
7 & 24t+264t^2+564t^3+264t^4+24t^5 \\
8 & 48t+933t^2+3513t^3+3513t^4+933t^5+48t^6 \\
\bottomrule
\end{array}
\]
Evidence: real-rooted for $k\leq 6$ and $n\leq 22$.  For $k=3$ and $k=6$
the rows also interlace consecutively in this range; for $k=4$ and $k=5$
the only failure is $D_{5,k}\interl D_{6,k}$, where
$D_{6,4}=D_{6,5}=D_{6,6}=9t+53t^2+47t^3+11t^4$ counts the $6$-cycles.

The case $k=2$, derangements by descents, is proved:
$D_{n,2}(t)$ has only real zeros for every $n$
\cite[Thm.~2.1]{Alexandersson2026Eulerian}, as conjectured in
\cite{FuLinZeng2018}; see the
\hyperref[derangementDescentPolynomial]{derangement page}.  Consecutive
interlacing is impossible there, since the decreasing involution gives
$\deg D_{2m,2}=2m-1$ while $\deg D_{2m-1,2}=2m-3$.

The polynomial $D_{n,3}(t)$ is palindromic.  Indeed, $\prod_{j\geq1}H[L_j]$
is $(1-\powerSum_1)^{-1}$, and $L_1=\completeH_1$, $L_2=\elementaryE_2$.
By Littlewood's identity $H\cdot H[\elementaryE_2]=\sum_\lambda\schurS_\lambda$
\cite[Ex.~I.5.4]{Macdonald1995}, so
\[
  \prod_{j\geq3}H[L_j]=\frac{(1-\powerSum_1)^{-1}}{\sum_\lambda\schurS_\lambda},
\]
which is fixed by $\omega$.  Since $\omega$ sends a descent set to its
complement (up to reversal), $\des\mapsto n-1-\des$ preserves the
distribution.  For $k=2$ and $k\geq4$ the corresponding quotient is not
$\omega$-invariant, and the data confirm that palindromicity fails; see the
\hyperref[realRootedCounterexamples]{counterexamples}.


\subsection[openProblemsEulerianInterlacing]{Interlacing of Eulerian-type polynomials}

\begin{conjecture}[Chow, \cite[Sec.~5]{Chow2022}]
\label{conjTypeDEulerianInterlacing}
Consecutive \hyperref[typeDEulerianPolynomial]{type $D$ Eulerian
polynomials} interlace: $D_n(t)\interl D_{n+1}(t)$ for $n\geq 2$.
\end{conjecture}

The polynomials $D_n(t)=\sum_{\sigma\in D_n}t^{\des_D(\sigma)}$ are
real-rooted, and their coefficients form \oeis{A066094}.  Stembridge's
identity $D_n(t)=B_n(t)-n2^{n-1}tA_{n-1}(t)$ relates them to the types $A$
and $B$.  The within-rank half interlacing of $A_{n-1}$, $B_n$ and $D_n$ is
proved \cite[Thm.~1]{Ma2026x}; see the
\hyperref[classicalTypeEulerianRealRooted]{classical types}.
For type $B$, consecutive interlacing follows from the recurrence
$B_n=((2n-1)t+1)B_{n-1}+2t(1-t)B'_{n-1}$.

\begin{conjecture}[Interlacing of binomial Eulerian polynomials]
\label{conjBinomialEulerianInterlacing}
The \hyperref[binomialEulerianPolynomial]{binomial Eulerian polynomials}
\[
  \tilde{A}_n(t)=\sum_{\sigma\in\symS_n}(1+t)^{\fix(\sigma)}t^{\exc(\sigma)}
\]
interlace consecutively: $\tilde{A}_n(t)\interl\tilde{A}_{n+1}(t)$.
\end{conjecture}

Real-rootedness of every $\tilde{A}_n(t)$, conjectured by
\name{Christos A. Athanasiadis}, is proved by \name{Jim Haglund} and
\name[P.B. Zhang]{Philip B. Zhang} \cite{HaglundZhang2019}.  Their
interlacing argument concerns a multivariate refinement and does not give
consecutive interlacing in $n$.  Computations support the conjecture.

The zig-zag Eulerian polynomials $Z_n(t)$, the $P$-Eulerian polynomials of
the zig-zag posets, were conjectured to be real-rooted by Petersen and
Zhuang \cite[Conj.~6.1(a)]{PetersenZhuang2024x}.\label{conjZigZagEulerian}
This is proved: $Z_n(t)/t$ has only simple negative zeros for $n\geq3$
\cite[Thm.~7.1]{Alexandersson2026Eulerian}, and independently, with strict
log-concavity, by Torres \cite[Thm.~1.1, Cor.~1.2]{Torres2026x}; see the
\hyperref[zigZagEulerianRealRooted]{permutation page}.  Petersen and Zhuang
also ask whether the generalizations $Z_n^r(t)$, which interpolate between
Eulerian and Narayana polynomials, are real-rooted; this remains open.

\begin{conjecture}[Gessel--Zhuang, \cite[Conj.~7.1]{GesselZhuang2024TwoSided}]
\label{conjInversePeaks}
For every $n\geq1$, the following polynomials are real-rooted:
\[
  \sum_{\substack{\pi\in\symS_n\\ \operatorname{pk}(\pi)=k}}
  t^{\operatorname{ipk}(\pi)+1},
  \qquad
  \sum_{\substack{\pi\in\symS_n\\ \operatorname{lpk}(\pi)=k}}
  t^{\operatorname{ilpk}(\pi)},
  \qquad
  \sum_{\substack{\pi\in\symS_n\\ \operatorname{lpk}(\pi)=k}}
  t^{\operatorname{ipk}(\pi)+1},
  \qquad
  \sum_{\substack{\pi\in\symS_n\\ \operatorname{pk}(\pi)=k}}
  t^{\operatorname{ilpk}(\pi)}.
\]
\end{conjecture}

Here $\operatorname{pk}$ counts peaks, $\operatorname{lpk}$ counts left
peaks (a peak at position $1$ is allowed when $\pi(1)\gt\pi(2)$), and
$\operatorname{ipk}(\pi)=\operatorname{pk}(\pi^{-1})$,
$\operatorname{ilpk}(\pi)=\operatorname{lpk}(\pi^{-1})$.  The conjecture
also covers the same inverse statistics over permutations with a fixed
number of up-down runs or biruns.  Gessel and Zhuang give
symmetric-function formulas for these two-sided distributions and verify the
conjecture for $n\leq50$.  Summing over $k$ gives the real-rooted
\hyperref[peakGeneratingPolynomial]{peak polynomials}, so the conjecture
refines a known result.  Their companion gamma-positivity conjecture is on
the \hyperref[gammaPositivityOpenProblems]{gamma-positivity page}.


\subsection[openProblemsBigDescents]{Big descents in pattern-avoiding permutations}

\begin{conjecture}[Elizalde--Rivera--Zhuang, \cite{ElizaldeRiveraZhuang2026}]
\label{conjBigDescents}
Let $\Pi\subseteq\symS_3$ consist of one or two patterns, and suppose that
$\Pi$ is not $\operatorname{bdes}$-Wilf equivalent to $\{123,132\}$.  Then
\[
  B_n(t;\Pi)\coloneqq\sum_{\pi\in\symS_n(\Pi)}t^{\operatorname{bdes}(\pi)}
\]
is real-rooted for every $n\geq0$.
\end{conjecture}

Here a \hyperref[bigDescent]{big descent} is an index $i$ with
$\pi_i\gt\pi_{i+1}+1$.  The case $\Pi=\{231\}$ is proved in
\cite{ElizaldeRiveraZhuang2026} via the equidistribution of big descents and
peaks on $231$-avoiders.  The two single-pattern cases that resist all
known methods are $132$ and $321$:
\[
\begin{array}{rll}
\toprule
n & B_n(t;132) & B_n(t;321) \\
\midrule
3 & 3+2t & 3+2t \\
4 & 4+9t+t^2 & 5+8t+t^2 \\
5 & 5+25t+12t^2 & 8+26t+8t^2 \\
6 & 6+55t+64t^2+7t^3 & 13+72t+45t^2+2t^3 \\
7 & 7+105t+233t^2+82t^3+2t^4 & 21+184t+196t^2+28t^3 \\
8 & 8+182t+674t^2+505t^3+61t^4 & 34+444t+732t^2+214t^3+6t^4 \\
\bottomrule
\end{array}
\]
With $s=1-t$, the generating functions are
\[
  \sum_{n\geq0}B_n(t;132)z^n
  =\frac{1-2sz+s(1-2t)z^2-\sqrt{\Delta}}{2tz(1-sz)},
  \quad
  \Delta=1-4z+6sz^2-4s^2z^3+s^2z^4,
\]
and
\[
  \sum_{n\geq0}B_n(t;321)z^n=\frac{2}{1-2sz^2+\sqrt{1-4z+4sz^2}}.
\]
Both families satisfy linear recurrences in $n$ of minimal order five, so no
three-term recurrence is available.  Evidence: real-rooted, with simple
zeros, and strictly interlacing consecutively for $n\leq40$ in both cases.
For $321$-avoiders, a
computer-assisted proof is formalized in Lean; a conceptual proof is still
missing.


\section[openProblemsWords]{Words, parking functions and Catalan words}

For a word $w=w_1\dotsm w_n$, let $\operatorname{pk}(w)$ count the indices
$1\lt i\lt n$ with $w_{i-1}\lt w_i\gt w_{i+1}$, and let $\exc(w)$ count
the indices $i$ with $w_i\gt i$.  Let $\mathrm{PF}_n$ be the set of
\hyperref[parkingFunctions]{parking functions} of length $n$.

\begin{conjecture}[Alexandersson, 2026]
\label{conjParkingExcedances}
The excedance polynomials
\[
  \sum_{w\in\mathrm{PF}_n}t^{\exc(w)}
\]
are real-rooted and interlace consecutively.
\end{conjecture}

\[
\begin{array}{rll}
\toprule
n & \sum_{w}t^{\exc(w)} & \sum_{w}t^{\operatorname{pk}(w)} \\
\midrule
2 & 2+t & 3 \\
3 & 6+9t+t^2 & 12+4t \\
4 & 24+71t+29t^2+t^3 & 59+66t \\
5 & 120+577t+515t^2+83t^3+t^4 & 351+825t+120t^2 \\
6 & 720+5023t+7808t^2+3029t^3+226t^4+t^5 & 2499+9704t+4604t^2 \\
7 & 5040+47349t+112588t^2+80589t^3+15975t^4+602t^5+t^6
  & 20823+115892t+115959t^2+9470t^3 \\
\bottomrule
\end{array}
\]
Evidence: real-rootedness and consecutive interlacing hold for $n\leq12$.
The peak polynomials, shown for comparison, are real-rooted only for
$n\leq9$; see the \hyperref[counterexamplesParkingPeaks]{counterexamples}.
Neither triangle is in the OEIS.  For comparison, the
descent polynomial of parking functions is
\hyperref[parkingFunctionDescentsRealRooted]{real-rooted}; it is $1/(n+1)$
times the descent polynomial of the words $[n+1]^n$.  Peaks behave
differently on words; see the
\hyperref[realRootedCounterexamples]{counterexamples}.

A \defin{Catalan word}\label{catalanWordDefinition} of length $n$ is a word
$w$ with $w_1=0$ and $0\leq w_{i+1}\leq w_i+1$.  These are the area
sequences of Dyck paths, that is, the increasing
parking functions shifted down by one.

\begin{conjecture}[Alexandersson, 2026]
\label{conjCatalanWordPeaks}
The peak polynomials of Catalan words,
\[
  C_n(t)\coloneqq\sum_{w}t^{\operatorname{pk}(w)}
  =\sum_{i\geq0}\frac{1}{n-i}\binom{n-i}{i}\binom{2n-3i}{n-2i-1}(t-1)^i,
\]
are real-rooted and interlace consecutively.
\end{conjecture}

The coefficients form \oeis{A116424}, which also counts Dyck paths of
semilength $n$ by occurrences of $UDUU$.  The generating function
$G=1+\sum_{n\geq1}C_n(t)x^n$ satisfies
\[
  G=1+x^2(1-t)G+x(1-x+tx)G^2,
\]
and \name{Jean-Luc Baril}, \name{Pamela E. Harris},
\name{Kimberly J. Harry}, \name{Matt McClinton}, and
\name{José L. Ramírez} give it in closed form
\cite[Thm.~18]{BarilHarrisHarryMcClintonRamirez2025}.
\[
\begin{array}{rl}
\toprule
n & C_n(t) \\
\midrule
3 & 4+t \\
4 & 9+5t \\
5 & 22+19t+t^2 \\
6 & 57+66t+9t^2 \\
7 & 154+221t+53t^2+t^3 \\
8 & 429+729t+258t^2+14t^3 \\
9 & 1223+2391t+1131t^2+116t^3+t^4 \\
\bottomrule
\end{array}
\]
Evidence: real-rooted and consecutively interlacing for $n\leq40$.

The conjecture has been proved by \name{ProgramComputer}
\cite[Thm.~1]{ProgramComputer2026CatalanPeaks}: every zero of $C_n(t)$ is
simple and negative, and consecutive polynomials strictly interlace.  The
proof makes an
\href{https://github.com/1337isnot1337/solved-conjectures/issues/59}{earlier
public argument} for all sufficiently large $n$ effective, with the explicit
bound $n\geq150$, and checks the remaining lengths exactly.  The project was
AI-assisted, and the preprint has not yet appeared on arXiv or in a journal.
A \href{https://github.com/ProgramComputer/open-math-research/blob/312293aa850010e9d9509d6850a482b9c7b7966f/topics/catalan-word-peaks/lean/CatalanStrictPeaksAllInterlacing.lean}{Lean
formalization} proves both statements for $C_n(t)$ defined by counting
Catalan words.
% Lean: CatalanStrictPeaks.strict_peaks_all_lengths; axioms propext,
% Classical.choice, Quot.sound in the repository's CI audit (not locally built).

By contrast, descents on Catalan words give
$\sum_k\binom{n-1}{2k}2^{n-1-2k}\mathrm{Cat}_kt^k$, a rescaled
gamma-polynomial of the
\hyperref[classicalNarayanaPolynomialsRealRooted]{Narayana polynomial},
which is therefore real-rooted.


\section[openProblemsSetPartitions]{Set partitions and trees}

\begin{conjecture}[Gessel--Lin--Zeng, \cite[Conj.~15]{GesselLinZeng2012}]
\label{conjJacobiStirling}
Every Jacobi--Stirling descent polynomial $A_{k,i}(t)$ is real-rooted.
\end{conjecture}

The polynomials $A_{k,i}(t)$ count Jacobi--Stirling permutations by
descents; the definition is in the
\hyperref[jacobiStirlingDescentRealRooted]{set-partition page}.  Each
summand $A_{k,S}(t)$ is real-rooted, and $A_{k,k}(t)$ is a second-order
Eulerian polynomial.  The conjecture is proved for
$i\in\{0,1,2,3,k-2,k-1,k\}$ \cite[Thm.~2]{MaWang2026x}, and remains open
for $4\leq i\leq k-3$.

\begin{conjecture}[Alexandersson, 2026]
\label{conjBigBlocks}
For $j\geq2$, let
\[
  P_{n,j}(t)\coloneqq\sum_{\pi}t^{\#\{\text{blocks of $\pi$ of size at least $j$}\}},
\]
summing over all \hyperref[setPartitions]{set partitions} $\pi$ of $[n]$.
Then $P_{n,j}(t)$ is real-rooted, and $P_{n,j}\interl P_{n+1,j}$.
\end{conjecture}

The exponential generating function is
\[
  \sum_{n\geq0}P_{n,j}(t)\frac{z^n}{n!}
  =\exp\Bigl(\sum_{r=1}^{j-1}\frac{z^r}{r!}
  +t\Bigl(e^z-\sum_{r=0}^{j-1}\frac{z^r}{r!}\Bigr)\Bigr).
\]
The triangles for $j=2$ and $j=3$ are \oeis{A124324} and \oeis{A355144}.
For $j=4$,
\[
\begin{array}{rl}
\toprule
n & P_{n,4}(t) \\
\midrule
4 & 14+t \\
5 & 46+6t \\
6 & 166+37t \\
7 & 652+225t \\
8 & 2780+1325t+35t^2 \\
9 & 12644+8062t+441t^2 \\
\bottomrule
\end{array}
\]
The polynomials $P_{n,j}$ were studied by \name{Jack Zhan} \cite{Zhan2023}.
The case $j=2$ is known: after a shift it is the descent polynomial of run-sorted
permutations \cite{AlexanderssonNabawanda2021,BonaHezo2016}.  The case $j=3$ is proved
(Alexandersson, 2026), so the first open case is $j=4$.  Two restrictions are proved for
every $j$, with consecutive interlacing (Alexandersson, 2026): partitions all of whose blocks
have size at least $j$, and partitions whose blocks have size $1$ or at least $j$, weighted
by the number of blocks of size at least $j$; see
\hyperref[bigBlockSetPartitions]{big blocks in set partitions}.  Evidence: $P_{n,j}$ is real-rooted and
interlaces consecutively for $j\leq6$ and $n\leq40$.  The analogous cycle statistic is the
\hyperref[conjLongCycleExcedances]{excedance conjecture above}.

\begin{conjecture}[Alexandersson, 2026]
\label{conjOrderedBigBlocks}
For $j\geq2$, the polynomials
\[
  O_{n,j}(t)\coloneqq\sum_{\pi}t^{\#\{\text{blocks of $\pi$ of size at least $j$}\}},
\]
summing over all ordered set partitions $\pi$ of $[n]$, are real-rooted, and
$O_{n,j}\interl O_{n+1,j}$.
\end{conjecture}

The exponential generating function is $1/(1-B_j(z,t))$, where
$B_j(z,t)=\sum_{r=1}^{j-1}z^r/r!+t\bigl(e^z-\sum_{r=0}^{j-1}z^r/r!\bigr)$ is the exponent
above.  For $j=3$ the first rows are $1$, $3$, $12+t$, $66+9t$, $450+91t$,
$3690+973t+20t^2$.  Evidence: real-rooted and consecutively interlacing for $j\leq5$ and
$n\leq30$.

\begin{conjecture}[Lin--Ma, \cite[Conj.~5.1]{LinMa2025}]
\label{conjReducedSchett}
For a multiset $M=\{1^{p_1},\dotsc,n^{p_n}\}$, let $\mathcal{T}_M$ be the
weakly increasing plane trees labeled by $M\cup\{0\}$, and let
$\operatorname{ee}(T)$ and $\operatorname{oe}(T)$ count the nodes with an
even number of children on even and odd levels, respectively.  Then every
\[
  \widehat{S}_{M,i}(t)\coloneqq
  \sum_{\substack{T\in\mathcal{T}_M\\
  \lfloor\operatorname{ee}(T)/2\rfloor=i}}
  t^{\lfloor\operatorname{oe}(T)/2\rfloor}
\]
is real-rooted.
\end{conjecture}

In a weakly increasing tree, labels weakly increase along every path from
the root, and the labels of the children of a node weakly increase from
left to right.  These polynomials refine the multiset Schett polynomials.
\name{Zhicong Lin} and \name{Jun Ma} prove their gamma-positivity and verify
the conjecture for plane trees and increasing trees with fewer than eleven
nodes.


\section[openProblemsGraphs]{Graphs, matroids and arrays}

\begin{conjecture}[Alexandersson, 2026]
\label{conjNestedMatroidZ}
Let $M$ be a \defin{nested matroid}, that is, a
\hyperref[latticePathMatroidDefinition]{lattice path matroid} of Ferrers shape:
for $1\leq q_1\leq\dotsb\leq q_r$, the bases are the sets
$\{b_1<\dotsb<b_r\}$ with $b_i\leq q_i$.  Then the $Z$-polynomial $Z_M(t)$ of
$M$ \cite{ProudfootXuYoung2018} is real-rooted.
\end{conjecture}

Real-rootedness of $Z_M(t)$ was conjectured for all matroids by
\name{Nicholas Proudfoot}, \name{Yuan Xu} and \name{Ben Young}
\cite{ProudfootXuYoung2018}, but \name{Jens Niklas Eberhardt} showed that it
fails for the extended binary Golay matroid \cite{Eberhardt2026x}; see the
\hyperref[realRootedCounterexamples]{counterexamples}.  It holds for sparse
paving matroids, including all uniform matroids, by \name{Matthew H. Y. Xie} and
\name{Philip B. Zhang} \cite[Thm.~1.1]{XieZhang2026SparsePaving}; in particular
rectangular shapes are covered.  $\gamma$-positivity of $Z_M(t)$ holds for every
matroid \cite[Thm.~1.8]{FerroniMatherneStevensVecchi2022x}.  The conjecture
above was checked exactly for all nested matroids with at most $24$ elements,
and for all nested matroids of rank at most $20$ with any number of elements.
For the shifted staircase $S_r$ ($q_i=2i$), the Catalan matroid satisfies
$Z(\mathrm{Cat}_{r+1})=(1+t)Z(S_r)$.  \name{Yiming Chen}, \name{Yao Li} and
\name{Ming Yao} express the $Z$-polynomials of all $(a,b)$-Catalan matroids
through uniform matroids \cite[Cor.~4.13]{ChenLiYao2025x}; for $S_r$ this gives
$rZ(S_r)=\sum_{m=1}^{r}Z(U_{m,2m})Z(S_{r-m})$.  The first values are
\[
Z(S_2)=1+3t+t^2,\quad Z(S_3)=1+8t+8t^2+t^3,\quad
Z(S_4)=1+22t+50t^2+22t^3+t^4,
\]
with $\gamma$-vectors $(1,1)$, $(1,5)$, $(1,18,8)$, $(1,59,91)$,
$(1,190,658,143)$, $(1,618,3899,2684)$, $(1,2047,20726,29812,3723)$ for
$r=2,\dotsc,8$.  Here $\gamma_1=\sum_{k=1}^{r}C_k-r$, with $C_k$ the Catalan
numbers.  Consecutive $Z(S_r)$ interlace for $r\leq100$.

\begin{conjecture}[Alexandersson, 2026]
\label{conjRibbonMatroidZ}
Let $M$ be a connected lattice path matroid of rank $r$ whose shape is a ribbon
(a connected skew shape without a $2\times2$ square, determined by the
composition of its row lengths).  Then $Z_M(t)$ equals the
Narayana polynomial (see \hyperref[narayanaNumberCSP]{Narayana numbers})
\[
N_{r+1}(t)=\sum_{k=0}^{r}\frac{1}{r+1}\binom{r+1}{k}\binom{r+1}{k+1}t^k,
\]
independently of the composition.  In particular it is real-rooted.
\end{conjecture}

This holds for the circuit $U_{r,r+1}$ \cite{ProudfootXuYoung2018} and was
checked for all $1022$ connected ribbon lattice path matroids on at most $11$
elements, including the fan matroids, whose $Z$-polynomials were computed by
\name{Linyuan Lu}, \name{Matthew H. Y. Xie} and \name{Arthur L. B. Yang}
\cite{LuXieYang2018x}.  These matroids are graphic.  A ribbon with loops or
coloops contributes a factor $(1+t)$ for each coloop; for instance, a hook
shape with $h$ nonempty rows gives $(1+t)^{r-h}N_{h+1}(t)$.  What explains this
rigidity?

\begin{conjecture}[Binder--Vecchi, \cite[Conj.~1.5]{BinderVecchi2026x}]
\label{conjRefinedChow}
For every loopless matroid, every refinement by cohomological degree of the
Hodge--Poincaré polynomial of the cohomology ring of its Bergman fan is
real-rooted.
\end{conjecture}

\name{Kyle Binder} and \name{Lorenzo Vecchi} prove this for uniform matroids
\cite[Thm.~1.4]{BinderVecchi2026x}, and the analogous statement for their
modified augmented fan \cite[Thm.~1.6]{BinderVecchi2026x}.  For the original
augmented fan the refinements need not be unimodal.  See the
\hyperref[matroidChowPolynomials]{Chow polynomials of matroids}.  The
unrefined Chow polynomials are not real-rooted in general; see the
\hyperref[realRootedCounterexamples]{counterexamples}.

\begin{conjecture}[Caiolo--Ferroni--Hoster, \cite[Conj.~7.3]{CaioloFerroniHoster2026x}]
\label{conjDualChow}
For every matroid $M$, the dual Chow polynomial $H^*_M(x)$ of its lattice
of flats has only real nonpositive zeros.  Moreover, the dual right
augmented Chow polynomial $F^*_M(x)$ is real-rooted, and its zeros are
interlaced by those of $H^*_M(x)$.
\end{conjecture}

The dual Chow polynomial is the Chow polynomial of the sign-twisted reverse
of the characteristic kernel; for a matroid of rank $r$, it is palindromic
of degree $r-1$ and gamma-positive.  For the Boolean matroid it is the
Eulerian polynomial, and for uniform matroids
\[
  H^*_{U_{r,n}}(x)=\binom{n-1}{r-1}
  +\sum_{j=0}^{r-2}\binom{n}{j}\binom{n-j-1}{r-j-1}
  A_j(x)\,(x+x^2+\dotsb+x^{r-j-1}).
\]
Both parts are proved for uniform matroids
\cite[Cor.~6.11]{CaioloFerroniHoster2026x}.  The analogous statement fails
for Cohen--Macaulay posets: the dual Chow polynomial
$4x^5+39x^4+120x^3+120x^2+39x+4$ of an EL-shellable poset has four
non-real zeros \cite[Rem.~7.4]{CaioloFerroniHoster2026x}.

\begin{conjecture}[Balanced arrays of uniform matroids]
\label{conjBalancedArrays}
For $1\leq r\leq d$ and $n\geq0$, let $a^{(r,d)}_{n,k}$ count the
$n\times d$ zero-one matrices in which every column sum is $k$ and every row
sum is at most $r$.  Then
\[
  P^{(r,d)}_n(q)\coloneqq\sum_{k\geq0}a^{(r,d)}_{n,k}q^k
  =\sum_{k\geq0}q^k\,[x_1^k\dotsm x_d^k]\,
  \bigl(1+\elementaryE_1+\dotsb+\elementaryE_r\bigr)^n
\]
is real-rooted for all $n$, $r$ and $d$.
\end{conjecture}

Here $\elementaryE_j$ is the elementary symmetric polynomial in
$x_1,\dotsc,x_d$; the possible rows are the independent sets of the uniform
matroid $U_{r,d}$.  The degree is $\lfloor rn/d\rfloor$.  The first open
case is $r=2$, $d=4$:
\[
\begin{array}{rl}
\toprule
n & P^{(2,4)}_n(q) \\
\midrule
2 & 1+6q \\
3 & 1+54q \\
4 & 1+204q+90q^2 \\
5 & 1+540q+2490q^2 \\
6 & 1+1170q+20610q^2+1860q^3 \\
7 & 1+2226q+102690q^2+106260q^3 \\
8 & 1+3864q+380100q^2+1543920q^3+44730q^4 \\
\bottomrule
\end{array}
\]
Evidence: for $r=2$, all $4\leq d\leq12$ and $n\leq40$ are real-rooted,
with consecutive interlacing in $n$.  For rank two there is a formula through matchings in complete
multipartite graphs, but real-rootedness of the matching polynomials does
not transfer to $q$.


\section[openProblemsPolytopes]{Tableaux, polytopes and Ehrhart theory}

For a standard Young tableau $T$, say that $i$ is a \emph{natural descent}
if $i+1$ lies in a strictly higher row than $i$, and let $\operatorname{ndes}(T)$
be the number of natural descents.

\begin{conjecture}[Alexandersson, 2026]
\label{conjSYTNaturalDescents}
For every partition $\lambda$, the natural-descent polynomial
\[
  N_\lambda(t)\coloneqq\sum_{T\in\SYT(\lambda)}t^{\operatorname{ndes}(T)}
\]
is real-rooted.
\end{conjecture}

Reading the cells of $\lambda$ row by row gives a natural labeling of the
Ferrers poset, and the natural descents of $T$ are the descents of the
corresponding linear extension.  Thus $N_\lambda(t)$ is the $P$-Eulerian
polynomial of the naturally labeled Ferrers poset, that is, the
\hyperref[hStarPolynomial]{$h^*$-polynomial} of its order polytope:
\[
  \sum_{m\geq0}\Omega_\lambda(m+1)\,t^m=\frac{N_\lambda(t)}{(1-t)^{|\lambda|+1}},
\]
where $\Omega_\lambda(m)$ counts the order-preserving maps from the Ferrers
poset to $[m]$.  The conjecture is the
\hyperref[neggersStanleyConjecture]{Neggers--Stanley} statement for Ferrers
posets.  The ordinary descents, where $i+1$ lies in a strictly lower row,
correspond to the column-strict labeling, and their polynomial is
real-rooted by \name{Francesco Brenti}'s theorem; see
\hyperref[sytDescentRealRooted]{descent polynomials of standard Young
tableaux}.  The polynomial $N_\lambda$ is palindromic exactly when all
removable corners of $\lambda$ lie on one antidiagonal; then the Ferrers
poset is graded and its Hibi ring is Gorenstein.
\[
\begin{array}{rl}
\toprule
\lambda & N_\lambda(t) \\
\midrule
(3,1) & 1+2t \\
(4,2) & 1+5t+3t^2 \\
(2,2,2) & 1+3t+t^2 \\
(3,2,1) & 1+7t+7t^2+t^3 \\
(4,2,2) & 1+13t+28t^2+13t^3+t^4 \\
(4,4,4) & 1+22t+113t^2+190t^3+113t^4+22t^5+t^6 \\
\bottomrule
\end{array}
\]
Rectangles are covered by Brenti's theorem on Gaussian posets
\cite{Brenti1989}.  Evidence: $N_\lambda$ is real-rooted for all $81{,}155$
partitions with $|\lambda|\leq35$.  The analogous statements for skew shapes,
for both natural and ordinary descents, are open as well.

\begin{conjecture}[Skew Schur Ehrhart polynomials]
\label{conjSkewSchurHStar}
For skew shapes $\lambda/\mu$ and $k\geq1$, the
\hyperref[hStarPolynomial]{$h^*$-polynomial} of
$n\mapsto|\SSYT(n\lambda/n\mu,k)|$ is real-rooted.
\end{conjecture}

For straight shapes, the counting function is a product of linear factors
with nonpositive roots, and real-rootedness of the $h^*$-polynomial is
due to \name{Francesco Brenti} \cite{Brenti1989}; see
\hyperref[schurEhrhartRealRooted]{Schur Ehrhart polynomials}.  For skew
shapes, the counting function itself can have non-real zeros: for
$\lambda/\mu=(3,2)/(1)$ and $k=3$ it is
$\frac14(n+1)^2(7n^2+10n+4)$, while its $h^*$-polynomial
$1+16t+22t^2+3t^3$ is real-rooted.  Evidence: all skew shapes with
$|\lambda|\leq8$ and $k\leq4$.

\begin{conjecture}[Alexandersson, May 2020]
\label{keyHStarRealRootedConjecture}\label{keyRealRootedHstar}
Let $\lambda$ be a partition and $\sigma$ a permutation, and let
$P_{\lambda,\sigma}(k)\coloneqq\key_{k\lambda,\sigma}(1^n)$, the
specialization of a \hyperref[key]{key polynomial}.  Define
$H_{\lambda,\sigma}(t)$ by
\[
  \sum_{k\geq0}P_{\lambda,\sigma}(k)t^k
  =\frac{H_{\lambda,\sigma}(t)}{(1-t)^{d+1}},
\]
where $d$ is the degree of $P_{\lambda,\sigma}$.  Then
$H_{\lambda,\sigma}(t)$ is real-rooted.
\end{conjecture}

Here $\key_{\lambda,\sigma}=\pi_\sigma(x^\lambda)$ is the
\hyperref[keyPolynomialOperatorDefinition]{key polynomial} indexed by a
partition and a permutation.  The value $\key_{k\lambda,\sigma}(1^n)$ is a
polynomial in $\lambda_1,\dotsc,\lambda_n$ and $k$, and for fixed $\lambda$,
$P_{\lambda,\sigma}$ is the Ehrhart polynomial of a union of faces of a
\hyperref[gtpolytopes]{Gelfand--Tsetlin polytope}
\cite{AlexanderssonAlhajjar2018}; see also
\hyperref[keyKogan]{key polynomials from Kogan faces}.  The same conjecture is stated
by \name{Alice L. L. Gao}, \name{Xin-Bei Liu}, \name{Arthur L. B. Yang},
and \name{James J. Y. Zhao} as the real-rootedness of the numerators of the
Lascoux series \cite[Conj.~6.3]{GaoLiuYangZhao2026x}.  Known cases, all for $\lambda=\delta_n$, include
Young-subgroup longest elements, where the numerator is an Eulerian
polynomial, and the permutations $[2,3,\dotsc,n,1]$ and
$[n,1,2,\dotsc,n-1]$, where it is a parking-function descent polynomial
\cite[Thm.~1.1]{GaoLiuYangZhao2026x}.  The
\hyperref[pitmanStanleyPolytope]{Pitman--Stanley polytopes} give a further
case.

A stronger form asked that the numerators $H_{\delta_n,\sigma}$ along every
saturated chain in Bruhat order form a Sturm sequence
\cite[Conj.~6.4]{GaoLiuYangZhao2026x}.  This fails: strong Bruhat covers
need not interlace, already in $\symS_6$.  Right weak-order covers
interlace for $n\leq7$ but not in general (Alexandersson, 2026); whether
all left weak-order covers interlace is open.

\begin{conjecture}[Hill--Luo--Trinh--Vindas-Meléndez,
\cite[Conj.~8.1]{HillLuoTrinhVindasMelendez2026x}]
\label{conjParkingPolytopeMagic}
For every $n\geq3$ and every $\mathbf{b}\in\setN_{\gt0}^n$, the
$\mathbf{b}$-parking-function polytope $\mathfrak{X}_n(\mathbf{b})$ is magic
positive.  In particular, its $h^*$-polynomial is real-rooted.
\end{conjecture}

The polytope $\mathfrak{X}_n(\mathbf{b})$ is the convex hull of the
$\mathbf{b}$-parking functions of length $n$.  A lattice polytope of
dimension $d$ is magic positive if its Ehrhart polynomial is
$\sum_i\mu_it^i(t+1)^{d-i}$ with all $\mu_i\geq0$; by a theorem of
\name{Petter Brändén}, this implies that the $h^*$-polynomial is
real-rooted.  The conjecture is proved for $\mathbf{b}=(a,b,\dotsc,b)$,
the only exception being $\mathbf{b}=(1,1)$
\cite[Thm.~7.2 and Cor.~7.12]{HillLuoTrinhVindasMelendez2026x}.  For
example, the $h^*$-polynomial of $\mathfrak{X}_3(3,2,2)$ is
$1+268t+661t^2+102t^3$.  Evidence: all $\mathbf{b}\in\{1,2,3,4\}^3$ and
several vectors of length four.


For margins $r\in\setN_{\gt0}^m$ and $c\in\setN_{\gt0}^n$ with equal sums,
the \defin{transportation polytope}\label{transportationPolytope}
$T(r,c)$ consists of all nonnegative real $m\times n$ matrices with row sums
$r$ and column sums $c$; for example,
$\mathcal{B}_n=T(1^n,1^n)$ is the
\hyperref[birkhoffPolytopeDefinition]{Birkhoff polytope}.

\begin{conjecture}[Alexandersson, 2026]
\label{conjTransportationHStar}
If both $r$ and $c$ have an entry at least $2$, then the
\hyperref[hStarPolynomial]{$h^*$-polynomial} of $T(r,c)$ is real-rooted.
\end{conjecture}

The hypothesis is needed: $\mathcal{B}_3$ has $h^*$-polynomial $1+t+t^2$,
and $\mathcal{B}_4$, $\mathcal{B}_5$ and a few other $T(r,1^n)$ fail too.
For example, $T((2,2,2),(2,2,2))=2\mathcal{B}_3$ has $h^*$-polynomial
$1+16t+25t^2+6t^3$.  Athanasiadis showed that the $h^*$-polynomial of
$\mathcal{B}_n$ is unimodal \cite{Athanasiadis2005}.  Evidence: all
margins up to size $5\times5$ with small entries.


\subsection[semiorderEulerianZeros]{Descents of linear extensions of unit interval orders}

Let $P$ be a \hyperref[unitIntervalOrder]{unit interval order} (a semiorder).
Label the unit intervals $1,\ldots,n$ by increasing left endpoint, with
$i\lt_P j$ when interval $i$ ends strictly before interval $j$ begins.
For a linear extension $w=w_1\cdots w_n$, count ordinary label descents,
$\des(w)=\#\{i:w_i\gt w_{i+1}\}$.  The resulting
\hyperref[pEulerianPolynomial]{$P$-Eulerian polynomial} is
\[
  W_P(t)=\sum_{w\in\mathcal{L}(P)}t^{\des(w)}.
\]

\begin{problem}[Unit interval orders]
Does $W_P(t)$ have only real zeros for every unit interval order $P$?
More strongly, suppose $P'$ is obtained by removing a single
comparability from $P$ and both posets admit the same left-endpoint
ordering.  Does $(1-\lambda)W_P(t)+\lambda W_{P'}(t)$ have only real
zeros for every $0\leq\lambda\leq1$?
\end{problem}

This is the \hyperref[neggersStanleyConjecture]{Neggers--Stanley question}
restricted to semiorders.  To compute $W_P$, let $G$ be the
incomparability graph of $P$ and partition its vertices into cliques.
Each partition contributes the product of its block weights: a singleton
has weight $1$, and a block of size $s\geq2$ has weight $tA_{s-1}(t)$,
where $A_j(t)$ is the \hyperref[eulerianPolynomial]{Eulerian polynomial}.
Indeed, splitting an extension before each left-to-right maximum gives
these blocks.  To recover the extension, order the blocks by increasing
maximum and put the maximum first in each block, ordering its tail freely.

In width at most two, $G$ is a disjoint union of paths.  Its nonsingleton
blocks form a matching, so $W_P(t)$ has only real zeros by the
\hyperref[matchingPolynomial]{matching polynomial theorem}
\cite{HeilmannLieb1972}.  For paths on three, four and five vertices this
gives $1+2t$, $1+3t+t^2$ and $1+4t+3t^2$, respectively.

Exact computation verifies the first statement through $n=12$,
including all $208{,}012$ Dyck paths at $n=12$.  For the stronger
question, exact root isolation gives a common interlacer for each of the
$101{,}763$ single-comparability changes through $n=10$.  This verifies
the whole interval $0\leq\lambda\leq1$ in that range.

\section[openProblemsGeometry]{Moduli spaces and Hurwitz numbers}

\begin{conjecture}[Ferroni--Pagaria--Vecchi, \cite[Conj.~4.13]{FerroniPagariaVecchi2026x}]
\label{conjTypeBModuli}
The Poincaré polynomials $B_n(x)$ of the type $B$ analogues
$\overline{\mathcal{M}}^B_{0,n}$ of $\overline{\mathcal{M}}_{0,n+1}$ are
real-rooted for all $n\geq1$.
\end{conjecture}

Here $\overline{\mathcal{M}}^B_{0,n}$ is the wonderful compactification of
the type $B$ braid arrangement with the minimal building set.  With $A_n(x)$
the Poincaré polynomial of $\overline{\mathcal{M}}_{0,n+1}$ and $B_0=1$,
\[
  B_{n+1}=B_n+x\sum_{j=1}^n\binom nj B_jB_{n-j}
  +x\sum_{j=2}^n\binom nj 2^jA_jB_{n+1-j}
\]
\cite[Thm.~4.11]{FerroniPagariaVecchi2026x}.  The first polynomials are
\[
  1,\quad 1+x,\quad 1+8x+x^2,\quad 1+35x+35x^2+x^3,\quad
  1+122x+482x^2+122x^3+x^4.
\]
Every $B_n$ is palindromic and gamma-positive, and real-rootedness is
verified for $n\leq300$.  This is the type $B$ analogue of the conjecture
of \name{Paolo Aluffi}, \name{Stephanie Chen} and \name{Matilde Marcolli}
that the polynomials $A_n(x)$ are real-rooted.

\begin{conjecture}[Coulter--Do--Moskovsky, \cite[Conjs.~3.13--3.14]{CoulterDoMoskovsky2023x}]
\label{conjMonotoneHurwitz}
For all $g\geq0$, $n\geq1$ and $\mu_1,\dotsc,\mu_n\geq1$, the $t$-monotone
Hurwitz number $H^t_{g,n}(\mu_1,\dotsc,\mu_n)$ is real-rooted in $t$, and
it interlaces each of the $n$ polynomials obtained by increasing one part
$\mu_i$ by one.
\end{conjecture}

Here a monotone factorization $(a_1\,b_1)\dotsm(a_r\,b_r)$ is weighted by
$t$ to the number of distinct values among $b_1,\dotsc,b_r$, and
$H^t_{g,n}(\mu)$ is the weighted count of transitive monotone factorizations
of length $|\mu|+2g-2+n$ of a permutation of cycle type $\mu$, divided by
$\mu_1\dotsm\mu_n$.  At $t=1$ these are the monotone Hurwitz numbers.  For
$(g,n)=(0,1)$ one obtains the \hyperref[classicalNarayanaPolynomialsRealRooted]{Narayana
polynomials}: $(d+1)H^t_{0,1}(d+1)$ is the Narayana polynomial of order
$d$.  Both parts are proved for $(g,n)=(0,1)$ and $(1,1)$ from a one-point
recursion.  The interlacing statement is checked in 1430 cases with
$g\leq5$, $n\leq5$ and $|\mu|\leq15$.

\name{Xavier Coulter} and \name{Norman Do} conjecture the same for the
two-parameter $bt$-monotone Hurwitz numbers, for every real
$b\notin\{0,-1\}$ \cite[Conjs.~4.21--4.22]{CoulterDo2025x}.  Here $b=0$
recovers the polynomials above, and $b=1$ is the orthogonal case.  The
evidence is more than 4000 checks over nine integer values of $b$.


\section[openProblemsTransformations]{Matrices and recurrences}

\begin{conjecture}[Mao--Wang, \cite[Conj.~4.1]{MaoWang2026x}]
\label{conj:eulerianDelannoySquares}
Let $A$ be the \hyperref[eulerianPolynomial]{Eulerian triangle} and let $D$
be the Delannoy triangle.
The row-generating functions of $A^2$ and $D^2$ have only real
nonpositive roots.
\end{conjecture}

The \hyperref[rowGeneratingFunction]{row-generating function} of a lower
triangular matrix $M$ is $M_n(x)=\sum_kM(n,k)x^k$.  \name{Jianxi Mao} and
\name{Lijie Wang} prove that for a broad class of matrices $B$,
real-rootedness of the rows of $M$ is inherited by the rows of $MB^r$
\cite[Thm.~1.3]{MaoWang2026x}; see the
\hyperref[narayanaTransformations]{Narayana transformations}.

\subsection[openProblemsElsewhere]{Related conjectures on other pages}

The \hyperref[jackSchurExpansionConjecture]{Schur expansion conjecture} for
Jack polynomials \cite{AlexanderssonHaglundWang2018} includes the
real-rootedness of the Jack rook and hit polynomials; a table and data for
partitions of size at most $9$ are given there.


\section[realRootedCounterexamples]{Counterexamples}

These families look real-rooted for small parameters, or were conjectured
to be real-rooted, but are not.

\subsection[counterexamplesWordPeaks]{Peaks on words}

Peaks on permutations give the real-rooted
\hyperref[peakGeneratingPolynomial]{peak polynomials}, and peaks on
\hyperref[conjCatalanWordPeaks]{Catalan words} are real-rooted.  Peaks on unrestricted words and on
parking functions fail.  For words over the alphabet $\{0,1\}$, the first five
polynomials $2$, $4$, $7+t$, $12+4t$, $21+10t+t^2$ are real-rooted, but
\[
  \sum_{w\in\{0,1\}^6}t^{\operatorname{pk}(w)}=37+22t+5t^2
\]
has discriminant $-256$.  Over the alphabets $[3]$, $[4]$ and $[5]$, the
first failures occur at $n=6$, $7$ and $7$, for example
$\sum_{w\in[3]^6}t^{\operatorname{pk}(w)}=292+334t+103t^2$.  Over the
alphabet $[n]$ itself, the polynomials are real-rooted for $n\leq7$, but
\[
  \sum_{w\in[8]^8}t^{\operatorname{pk}(w)}
  =1215616+5640832t+7878016t^2+2042752t^3
\]
has two non-real zeros, and so do the rows $n=9,10$.

\subsection[counterexamplesParkingPeaks]{Peaks on parking functions}

The peak polynomials $\sum_{w\in\mathrm{PF}_n}t^{\operatorname{pk}(w)}$ of
\hyperref[conjParkingExcedances]{parking functions} are real-rooted and
interlace consecutively for $n\leq9$, but
\[
  \sum_{w\in\mathrm{PF}_{10}}t^{\operatorname{pk}(w)}
  =24188458+278306304t+970166022t^2+930536936t^3+154749971t^4
\]
has two non-real zeros, and so do the rows $n=11,12$.  In contrast, the
excedance polynomials of parking functions remain real-rooted through
$n=12$.

\subsection[counterexamplesPalindromic]{Long-cycle descent polynomials}

The descent polynomials $D_{n,k}(t)$ of
\hyperref[conjLongCycleDescents]{permutations with all cycles of length at
least $k$} are palindromic for $k=1$ and $k=3$.  For other $k$, early rows
can be palindromic, but this does not persist: $D_{n,2}$ is not palindromic
for $n\geq4$, $D_{6,k}=9t+53t^2+47t^3+11t^4$ for $k=4,5,6$, and
$D_{n,4}$ is not palindromic for $10\leq n\leq 22$.

\subsection[counterexamplesPosets]{Posets and spheres}

The \hyperref[neggersStanleyConjecture]{Neggers--Stanley conjecture}, that
$P$-Eulerian polynomials of naturally labeled posets are real-rooted, is
false; counterexamples were constructed by \name{John R. Stembridge}
\cite{Stembridge2007NeggersStanleyCounter}.

The real root conjecture, posed by \name{Tadeusz Januszkiewicz} and
observed by \name{Victor Reiner} and \name{Volkmar Welker} to imply the
Neggers--Stanley conjecture for graded posets of width two, asserts that the
$h$-polynomial of a flag homology sphere is real-rooted.
\name{Świętosław Gal} shows that it fails in every dimension at least five
\cite{Gal2005}.  His smallest known example is a flag triangulation of the
$5$-sphere with $f$-polynomial
$1+17t+109t^2+345t^3+575t^4+483t^5+161t^6$, whose $h$-polynomial
$1+11t+39t^2+59t^3+39t^4+11t^5+t^6$ has only two real zeros.  Gal's
weaker \hyperref[galConjecture]{gamma-positivity conjecture} is open.

The antichain polynomials of the boxes $[2]\times[n]\times[k]$ are
real-rooted \cite{Jiang2026x}, but
$N_{[3]\times[3]\times[3]}(x)$ is not; see
\hyperref[antichainBoxStability]{antichain polynomials}.

\subsection[counterexamplesPetersen]{Generalized Petersen graphs}

\name{Rohan Pandey} conjectured that the
\hyperref[realRootedIndependence]{independence polynomial} of the
generalized Petersen graph $\mathrm{GP}(n,k)$, with $n\geq2k+1$, is
real-rooted if and only if $k$ is even
\cite[Conj.~4.1]{Pandey2026x}.\label{conjPetersenIndependence}
The evidence there is numerical.  The conjecture fails already for $k=2$:
\[
  I(\mathrm{GP}(9,2),x)
  =1+18x+126x^2+438x^3+801x^4+747x^5+303x^6+27x^7
\]
has only five real zeros; the remaining pair is approximately
$-0.857\pm0.056\,i$.  Also $I(\mathrm{GP}(10,2),x)$,
$I(\mathrm{GP}(11,2),x)$ and $I(\mathrm{GP}(11,4),x)$ are not real-rooted.
These counts were verified by exact Sturm sequences.

\subsection[counterexamplesTransformations]{Transformations, matroids and characters}

\name{Francesco Brenti} conjectured that the
\hyperref[eulerianTransformation]{Eulerian transformation}
$x^n\mapsto A_n(x)$ preserves real-rootedness.  This is false in general
\cite{BrandenJochemko2022}.

The general real-rootedness conjectures for chain polynomials, Chow
polynomials and augmented Chow polynomials of matroids are false: there are
matroids for which these polynomials are not even log-concave
\cite[Thms.~A--C]{BercziCoronEurLi2026x}.  The chain-polynomial examples
have rank four, and the Chow-polynomial examples have rank seven, both
minimal; see \hyperref[realRootsChow]{Chow polynomials}.

The $Z$-polynomials \cite{ProudfootXuYoung2018} and Kazhdan--Lusztig
polynomials of matroids were conjectured to be real-rooted.  Both fail for the
extended binary Golay matroid \cite{Eberhardt2026x}; see also the
\hyperref[conjNestedMatroidZ]{nested-matroid conjecture}.

\name{Christian Gaetz} and \name{Laura Pierson} conjectured that the
coefficient polynomials of identity-pattern character polynomials are
real-rooted with all roots below the pattern length
\cite[Conj.~1.4]{GaetzPierson2022x}.  \name{Jonas Iskander} gives a
degree-$11$ coefficient polynomial for the identity pattern of length $14$
with only nine real roots \cite[Ex.~1.1]{Iskander2023x}; see
\hyperref[characterSymmetricFunctions]{character polynomials}.

For skew shapes, the counting function
$n\mapsto|\SSYT(n\lambda/n\mu,k)|$ need not be real-rooted, although its
\hyperref[conjSkewSchurHStar]{$h^*$-polynomial appears to be}.
